\section{The model and the exact law}\label{sec:multistate-model}

Fix $d\ge2$ and a start law $\mu$ on $\{1,\ldots,d\}$. Let $(X_j)_{j\ge1}$ be the
Markov chain with $X_1\sim\mu$ that repeats its state with probability
$p\in[0,1]$ and otherwise moves to each of the other $d-1$ states with
probability $r=(1-p)/(d-1)$. Put $\sigma=p-r=(dp-1)/(d-1)$. The
occupation vector $K_N=(K_{N,1},\ldots,K_{N,d})$ counts the visits to each state
among $X_1,\ldots,X_N$, and $P^\mu_N(k)=\Prob(K_N=k)$. A word of length $N$ that
starts in $a$ and has $j$ changes has probability $\mu_ap^{N-1-j}r^j$. A bin $k$
with $\supp k=\{i:k_i>0\}=F$ lies on the face $F$ of the simplex. The vertices are
the bins $Ne_i$, with $V^\mu_i=P^\mu_N(Ne_i)=\mu_ip^{N-1}$. The uniform start
$\mu_i=1/d$ is the default in Sections~\ref{sec:simplex-thresholds},
\ref{sec:frontier-modes} and~\ref{sec:frontier-boundary}, and for it we write
$P_N(k)$ and $V_N=p^{N-1}/d$. With the uniform start and $0<p<1$, the chain is
also the one-dimensional $d$-state Potts model with free boundaries, in which a
configuration has weight proportional to $(p/r)^{\#\{s<N:\,x_s=x_{s+1}\}}$. For
$d=2$ this is the Ising chain of Paper~A~\cite{paperA}.

We use the notation shared by the 4 papers of this program. Let $L=\log N$ and
$T=Nr$, which is $N$ times the switching probability to one given other state.
Then the transition matrix is $P_N=I+(T/N)Q$ with $Q=J-dI$, where $J$ is the
all-ones matrix, so $Q$ switches at rate $1$ to each other state. We write
$\tau=T/L$. A face $F$ with $k=|F|$ states has exit rate $\lambda_F=d-k$, the
principal Dirichlet eigenvalue of $-Q$ on $F$. The value of $T$ at a crossing is
called $z_N$, with more indices when needed. Paper~C calls $K_F$ the
second-order constant of a face $F$. Here $K_F=\mu(F)k^{k-1/2}(4\pi)^{-(k-1)/2}$
with $\mu(F)=\sum_{a\in F}\mu_a$ for a face with $k\ge2$ states, and
$K_{\{i\}}=\mu_i$ for a vertex (Section~\ref{sec:general-start}). For $p>0$ we
also use $u=r/p$ and $v=r/\sigma=u/(1-u)$, and in proofs the variable $z=Nu=T/p$. Exact statements
are written in $u$ and asymptotic statements in $T$. Note that
$z-T=(d-1)Tz/N$. So an asymptotic statement whose error term contains $T^2/N$
also holds with $z$ in place of $T$. At finite $N$ the 2 differ, and
Remark~\ref{rem:K3-finite} says which one it uses.

For a positive count vector $n=(n_1,\ldots,n_k)$ with total $N$ and a state $a$,
let $W^{(a)}_n(j)$ be the number of words with counts $n$ and $j$ changes that
start with $a$, and let $W_n(j)=\sum_aW^{(a)}_n(j)$. Put
$S^{(a)}_n(u)=\sum_jW^{(a)}_n(j)u^j$ and $S_n(u)=\sum_aS^{(a)}_n(u)=\sum_jW_n(j)u^j$.
These polynomials depend only on the positive coordinates, and not on $d$.

\begin{proposition}[the generating function]\label{prop:dgf}
Put $y_i=x_it/(1-\sigma x_it)$ and $P^\mu_0(0,\ldots,0)=1$. As formal power
series,
\[
\sum_{N\ge0}\sum_{k_1+\cdots+k_d=N}P^\mu_N(k)\,x_1^{k_1}\cdots x_d^{k_d}\,t^N
=1+\frac{\sum_i\mu_iy_i}{1-r\sum_iy_i}.
\]
\end{proposition}

\begin{proof}
Let $F_i$ be the generating function of the words that end in state $i$, and
$S=\sum_iF_i$. Conditioning on the last letter gives
$F_i=x_it\bigl(\mu_i+pF_i+r(S-F_i)\bigr)$, that is, $F_i=y_i(\mu_i+rS)$. Summing
over $i$ gives $S=\sum_i\mu_iy_i+rS\sum_iy_i$, and the generating function is
$1+S$.
\end{proof}

\begin{proposition}[the exact law]\label{prop:dpmf}
Let $N\ge1$, let $k$ have nonnegative integer coordinates with sum $N$, and put
$F=\supp k$. Sums over $m$ run over $m_i=0$ for $i\notin F$ and $1\le m_i\le k_i$
for $i\in F$, and $M=\sum_im_i$.
\begin{enumerate}[(a)]
\item (refresh form) With $0^0=1$,
\[
P^\mu_N(k)=\sum_m\Bigl(\sum_{i\in F}\mu_im_i\Bigr)\frac{(M-1)!}{\prod_im_i!}\,
r^{M-1}\sigma^{N-M}\prod_{i\in F}\binom{k_i-1}{m_i-1}.
\]
For the uniform start this reads
\[
P_N(k)=\frac1d\sum_m\frac{M!}{\prod_im_i!}\,
r^{M-1}\sigma^{N-M}\prod_{i\in F}\binom{k_i-1}{m_i-1}.
\]
\item (run form) Let $p>0$ and let $n$ be the vector of positive coordinates of
$k$. Then $P^\mu_N(k)=p^{N-1}\sum_{a\in F}\mu_aS^{(a)}_n(u)$. In particular
$S_n(u)=P_N(k)/V_N$.
\item (ratio form) If $1/d<p<1$, so that $0<u<1$ and $v>0$, then
\begin{equation}\label{eq:simplex-positive-sum}
S_n(u)=(1-u)^{N-1}\sum_mM!\,v^{M-1}\prod_{i\in F}\frac1{m_i!}\binom{n_i-1}{m_i-1},
\end{equation}
which is a sum of positive terms.
\end{enumerate}
\end{proposition}

\begin{proof}
(a) By the proof of Proposition~\ref{prop:dgf},
$S=\sum_{M\ge1}r^{M-1}\bigl(\sum_i\mu_iy_i\bigr)\bigl(\sum_iy_i\bigr)^{M-1}$. By the
multinomial theorem the coefficient of $\prod_iy_i^{m_i}$ in
$(\sum_i\mu_iy_i)(\sum_iy_i)^{M-1}$ is $\sum_i\mu_im_i(M-1)!/\prod_jm_j!$. With
$z_i=x_it$ we have $y_i=z_i/(1-\sigma z_i)$, and
$[z^k](z/(1-\sigma z))^m=\binom{k-1}{m-1}\sigma^{k-m}$ for $k\ge m\ge1$. For
$m=0$ only $k=0$ contributes. Multiplying the coefficients proves (a), and for
the uniform start $\sum_i\mu_im_i=M/d$.
(b) A word that starts with $a$ and has $j$ changes has probability
$\mu_ap^{N-1}u^j$.
(c) Divide the uniform case of (a) by $V_N=p^{N-1}/d$, and note that
$r^{M-1}\sigma^{N-M}/p^{N-1}=(1-u)^{N-1}v^{M-1}$ since $r/p=u$ and $\sigma/p=1-u$.
\end{proof}

At $p=1/d$ we have $\sigma=0$ and only $M=N$ survives, which gives the
multinomial law for the uniform start. At $p=1$ only $M=1$ survives, which gives
mass $\mu_i$ at $Ne_i$. For $\sigma<0$ the terms in (a) can cancel, so for
numerical work we use exact arithmetic or dynamic programming.

\subsection{Bessel bounds}

For $x\ge0$ the modified Bessel function of order $1$ is
$I_1(x)=\sum_{j\ge0}(x/2)^{2j+1}/(j!(j+1)!)$, and comparing coefficients gives
$I_1(x)\le\sinh x\le e^x$. We also use the expansion~\cite[Eq.~10.40.1]{dlmf}
\[
I_1(x)=\frac{e^x}{\sqrt{2\pi x}}\Bigl(1-\frac3{8x}+O(x^{-2})\Bigr),
\qquad x\to\infty,
\]
and the function
\[
\Phi(x)=\sum_{j\ge0}\frac{x^j}{j!(j+1)!}=\frac{I_1(2\sqrt x)}{\sqrt x},
\qquad\Phi(0)=1 .
\]
Since $\binom{2j}j\le4^j$, we have $1/(j!(j+1)!)\le4^j/(2j)!$, and comparing term
by term with the even part of the series of $e^{2\sqrt x}$ gives
$\Phi(x)\le e^{2\sqrt x}$. Lemmas~\ref{lem:bessel-law}, \ref{lem:clipped-product}
and~\ref{lem:gaussian-phi-average} in Appendix~\ref{app:thresholds} give the facts
about $\Phi$ that we use, a bound for products, and an expansion of Gaussian
averages. The Laguerre--Bessel comparison needed near the boundary is
Lemma~\ref{lem:uniform-laguerre-bessel}.
